% Generated draft: requires mathematical and source review.
\documentclass[11pt]{article}
\usepackage[margin=25mm]{geometry}
\usepackage{fontspec}
\setmainfont{Latin Modern Roman}
\usepackage{amsmath,amssymb,mathtools,unicode-math}
\setmathfont{Latin Modern Math}
\usepackage{enumitem,needspace,xurl,hyperref}
\setlength{\parindent}{0pt}
\setlength{\parskip}{5pt}
\setlength{\emergencystretch}{3em}
\newcommand{\N}{\mathbb{N}}
\newcommand{\Z}{\mathbb{Z}}
\newcommand{\Q}{\mathbb{Q}}
\newcommand{\R}{\mathbb{R}}
\newcommand{\C}{\mathbb{C}}
\newcommand{\F}{\mathbb{F}}
\newcommand{\A}{\mathbb{A}}
\newcommand{\PP}{\mathbb{P}}
\newcommand{\E}{\mathbb{E}}
\newcommand{\eps}{\varepsilon}
\newcommand{\dd}{\,\mathrm{d}}
\begin{document}
\section*{UnsolvedMath ID: 30005258}
\textbf{Source problem code:} catalogue entry
\subsection{Definitions and mathematical statement}
\textbf{Source-attributed conjecture (Bourgain--Spencer).} Let \(A\) be a stationary random conductivity field in spatial dimension \(d\). Stationarity means that its probability law is invariant under translations of its spatial domain. Write \(\overline A(\nabla)\) for the convolution operator associated with the averaged solution operator, using the notation of the supplied report, and write \(\overline A(\xi)\) for its Fourier multiplier symbol. Thus, whenever this representation is defined,
\[
\widehat{\overline A(\nabla)f}(\xi)
 =\overline A(\xi)\widehat f(\xi).
\]
The supplied excerpt refers to an earlier construction of this averaged operator that is not included in the record. Consequently, this notation identifies the intended object but does not specify its construction or normalization completely.

The conjecture asks whether sufficiently strong mixing of \(A\) implies that \(\overline A\) is almost \(C^{2d}\) in a neighborhood of \(\xi=0\). Here mixing concerns decay of statistical dependence between spatially separated portions of the field. No particular mixing condition or quantitative decay rate is supplied.

For nonintegral \(s=k+\alpha>0\), with \(k\in\Z_{\geq0}\) and \(0<\alpha<1\), \(C^s\) denotes the local H\"older class: derivatives through order \(k\) exist and are continuous, and derivatives of order \(k\) are locally \(\alpha\)-H\"older continuous, componentwise for a matrix-valued symbol. A possible precise interpretation of the requested regularity is
\[
\forall\varepsilon\in(0,1)\quad
\exists r_\varepsilon>0:\qquad
\overline A\in C^{2d-\varepsilon}
       \bigl(B_{r_\varepsilon}(0)\bigr).
\]
This is an interpretation of ``almost \(C^{2d}\)'', not an endpoint assertion supplied by the source. The record does not determine whether a single neighborhood must work for every \(\varepsilon\), whether the mixing assumption may depend on \(\varepsilon\), or whether a different endpoint convention is intended. The conjecture therefore remains a formulation schema until these choices and the admissible conductivity fields are specified.

\textbf{Explicit discrete setting supplied as background.} The excerpt subsequently specializes its discussion of a perturbative result to independent, identically distributed, uniformly elliptic random conductivities \(A=(A_x)_{x\in\Z^d}\), with
\[
L=\nabla^*A\nabla\quad\text{on }L^2(\Z^d).
\]
Here \((\nabla f)_j(x)=f(x+e_j)-f(x)\), where \(e_j\) is the \(j\)-th coordinate vector, and \(\nabla^*\) is the adjoint for counting measure. In this lattice setting the multiplier is defined on the frequency torus, with regularity near zero understood in local coordinates. The record does not establish that this discrete specialization is the exclusive domain of the general conjecture.

The excerpt attributes to Bourgain and to Kim--Lemm the following perturbative statement: for \(d\geq2\),
\[
A_x=1+\delta B_x,\qquad
|B_x|\leq1,\qquad \E[B_x]=0,\qquad 0<\delta\ll1,
\]
one has
\[
\overline A(\nabla)=\operatorname{Id}+L_\delta,
\qquad
|L_\delta(x,y)|
 \leq C\delta^2\langle x-y\rangle^{C\delta-3d},
\qquad \langle z\rangle:=1+|z|,
\]
for a constant \(C>0\) described there as universal. The source describes the resulting symbol regularity as \(C^{2d-C\delta}\), and calls the kernel decay essentially optimal in the sense that it cannot be improved beyond \(\langle\cdot\rangle^{-3d}\). These are reported background claims, not hypotheses restricting the general conjecture or independently verified results here. The exact smallness threshold, parameter dependence of the constant, and interpretation at integral regularity exponents are not specified in the supplied excerpt.
\subsection{Short English statement}
The Bourgain--Spencer conjecture asks whether sufficiently strong mixing of a stationary random conductivity field forces the symbol of its averaged operator to have regularity arbitrarily close to \(C^{2d}\) near zero frequency. The supplied record leaves the mixing hypothesis, full averaged-operator construction, general spatial setting, and meaning of ``almost'' unspecified. Its dated curation labels the problem partially solved, citing perturbative nearly \(C^{2d}\) results and nonperturbative nearly \(C^d\) progress; this assessment is supplied rather than independently verified.
\subsection{Sources}
\begin{itemize}
\item Local UnsolvedMath registry, supplied source record with permanent ID \(30005258\), exact source problem code \texttt{catalogue entry}, title ``Bourgain--Spencer Regularity Conjecture'', in the Oberwolfach Reports Open Problems collection. The registry labels the statement \texttt{corrected\_verified} and states that it was reconstructed and checked against the official report; these are supplied provenance assertions, not verification performed for this draft.
\item Supplied report citation: \emph{Arbeitsgemeinschaft: Quantitative Stochastic Homogenization (2023)}, Oberwolfach Reports, pp.\ 2741--2809, DOI \texttt{10.4171/owr/2022/48}. \url{https://doi.org/10.4171/owr/2022/48}. The supplied literature list labels this ``Oberwolfach Report 48/2022'', and the excerpt includes the page marker 2783. Both year descriptions are retained without resolving their relationship.
\item Supplied literature reference: Duerinckx, \emph{Journal de Mathematiques Pures et Appliquees}, DOI \texttt{10.1016/j.matpur.2023.06.005}. \url{https://doi.org/10.1016/j.matpur.2023.06.005}. No article title or full bibliographic details are supplied.
\item Attribution in the excerpt: the conjecture is named for Bourgain--Spencer; the perturbative theorem is attributed to Bourgain [1] and Kim--Lemm [5]. The excerpt also cites [4] for more general results. The bibliography identifying these numbered references is not supplied. No separate proposer is listed in the registry.
\item Literature assessment date: \texttt{2026-08-22}, copied exactly from \texttt{literature\_checked\_at}. This is supplied dated curation, not independently verified. It reports nearly \(C^{2d}\) perturbative regularity and intermediate nearly \(C^d\) nonperturbative regularity, while describing the general strong-mixing conjecture as unresolved. The registry status is \texttt{partially\_solved}. The source explicitly requests independent verification of time-sensitive status and cited claims.
\end{itemize}
\end{document}
